\documentclass[11pt]{article}

\usepackage[utf8]{inputenc}
\usepackage[T1]{fontenc}
\usepackage{inconsolata}
\usepackage{geometry}
\geometry{a4paper, margin=0.85in}
\usepackage{xcolor}
\usepackage{tcolorbox}
\tcbuselibrary{skins, breakable}
\usepackage{enumitem}
\usepackage{listings}
\usepackage{caption}
\usepackage{hyperref}

% ── Listings style for prompt display ─────────────────────────────────────
\lstdefinestyle{promptstyle}{
  basicstyle=\small\ttfamily,
  backgroundcolor=\color{userbg},
  frame=single,
  rulecolor=\color{userheader},
  framesep=6pt,
  xleftmargin=2pt,
  xrightmargin=2pt,
  breaklines=true,
  showstringspaces=false,
  columns=fullflexible,
  keepspaces=true,
  escapeinside={(*@}{@*)},
}

% ── Colors ────────────────────────────────────────────────────────────────
\definecolor{promptgreen}{RGB}{0,120,50}
\definecolor{systembg}{RGB}{255,248,235}
\definecolor{systemheader}{RGB}{160,120,50}
\definecolor{userbg}{RGB}{237,247,237}
\definecolor{userheader}{RGB}{70,130,70}
\definecolor{developbg}{RGB}{245,245,245}
\definecolor{developheader}{RGB}{80,80,80}
\definecolor{avoidbg}{RGB}{255,243,243}
\definecolor{avoidheader}{RGB}{175,75,75}
\definecolor{fewshotbg}{RGB}{240,245,255}
\definecolor{fewshotheader}{RGB}{60,100,160}
\definecolor{codegray}{RGB}{245,245,245}

% ── Prompt box (no title bar, just a thin colored border) ─────────────────
\newtcolorbox{promptbox}[3][]{%
  enhanced, breakable,
  colback=#2, colframe=#3,
  sharp corners, boxrule=0.6pt,
  left=8pt, right=8pt, top=8pt, bottom=8pt,
  before upper={\parindent=0pt\parskip=6pt\small},
}

\newcommand{\pgreen}[1]{{\color{promptgreen}\texttt{#1}}}

\begin{document}

\appendix
\section{Prompts for Judge Program Generation}

\subsection{User Prompt for LLM-as-a-Judge}\label{app:llm-judge}

\begin{promptbox}{userbg}{userheader}
\textbf{User Prompt for LLM-as-a-Judge}

\medskip
You are an expert evaluator assessing AI-generated responses. Determine which of the two responses better serves the user's needs.

\pgreen{<question>}\pgreen{[question]}\pgreen{</question>}

\pgreen{<response\_a>}\pgreen{[response\_a]}\pgreen{</response\_a>}

\pgreen{<response\_b>}\pgreen{[response\_b]}\pgreen{</response\_b>}

Reply with ONLY ``A'' or ``B''.
\end{promptbox}

\subsection{Generation Prompt}\label{app:gen-prompt}

\begin{promptbox}{userbg}{userheader}
\textbf{Generation Prompt}

\medskip
You are an expert Python developer and AI evaluation researcher.

Your task: Write a Python function that evaluates the quality of an LLM-generated response to a given query. The function should return a numeric score where HIGHER values indicate BETTER quality.

\textbf{EVALUATION STRATEGY} --- focus STRICTLY on this dimension:

\pgreen{[Heuristic Name]}: \pgreen{[Heuristic Description]}

\pgreen{[Avoid-Repetition Block, if prior variants exist]}

Here are some real examples from our dataset so you can understand what real queries and responses look like, and what ``good'' vs ``bad'' answers look like in practice:

\pgreen{[Few-Shot Examples $\times$ 10]}

\textbf{IMPORTANT REQUIREMENTS:}

-- Output ONLY valid, executable Python code inside \texttt{```python ... ```} blocks. No explanation.

-- The function signature must be exactly: \texttt{def judging\_function(query, response):}

-- Return a single numeric score (int or float). Higher = better quality.

-- Use standard Python string/math operations for speed. You may use common libraries (re, math, collections, string, statistics) but NO heavy ML/NLP libraries.

-- This is variant \pgreen{[variant\_index]}/\pgreen{[programs\_per\_heuristic]} --- use a MEANINGFULLY DIFFERENT algorithm, features, and scoring formula than other variants.

-- Include comprehensive error handling (try/except) so the function never crashes.

-- The function must handle edge cases (empty strings, very short/long inputs).

-- Return scores in a reasonable numeric range (e.g., 0--10 or 0--100).

-- Make the function DISCRIMINATIVE: it should produce clearly different scores for high-quality vs low-quality responses.

\begin{verbatim}
```python
def judging_function(query, response):
    # Your implementation here
```
\end{verbatim}
\end{promptbox}

\subsection{Quality-Evaluation Heuristics}\label{app:heuristics}

\begin{promptbox}{systembg}{systemheader}
\textbf{Heuristic Definitions (10 Dimensions)}

\medskip
\textbf{1.~Relevance to the Query}\quad
Evaluate how semantically relevant the response is to the question asked. Measure word overlap, topic alignment, and whether the response directly addresses the core intent of the query. Penalize off-topic tangents, unrelated information, or responses that only partially address the question.

\textbf{2.~Language Quality and Readability}\quad
Evaluate language quality and readability. Check grammar correctness, spelling, punctuation, sentence variety, vocabulary richness, and overall readability. Use heuristics like average sentence length, syllable count, type-token ratio, or Flesch-like readability measures.

\textbf{3.~Completeness and Coverage}\quad
Evaluate completeness and thoroughness of the answer. Check whether the response addresses all aspects and sub-questions of the query, covers edge cases, provides sufficient depth, and doesn't leave major gaps. Penalize partial or superficial answers.

\textbf{4.~Factual Accuracy Indicators}\quad
Evaluate indicators of factual reliability. Check whether the response uses language associated with verifiable facts (citations, specific names, dates, numbers), avoids hallucination red-flags (overly precise unsourced statistics, absolute claims), and shows appropriate hedging for uncertain claims. Penalize sensationalism and conspiracy-style language.

\textbf{5.~Logical Coherence and Argument Structure}\quad
Evaluate logical coherence. Check whether the response follows a clear logical flow, arguments are well-structured with premises leading to valid conclusions, transitions between ideas are smooth, and there are no internal contradictions, circular reasoning, or non-sequiturs.

\textbf{6.~Clarity and Conciseness}\quad
Evaluate clarity and conciseness. Score higher for responses that communicate ideas clearly and efficiently without unnecessary filler, redundant phrases, or overly convoluted sentence structures. Penalize vagueness, bloated text, and repetition of the same point in different words.

\textbf{7.~Reasoning Transparency and Step-wise Formulation}\quad
Evaluate how transparently the response shows its reasoning process. Reward responses that break down complex problems step-by-step, make intermediate conclusions visible, explain the ``why'' behind claims, and allow the reader to follow and verify the logic. Penalize opaque answers that jump directly to conclusions without showing reasoning.

\textbf{8.~Epistemic Calibration and Uncertainty Communication}\quad
Evaluate how well the response communicates confidence and uncertainty. Reward responses that distinguish between well-established facts and speculative claims, use appropriate hedging language (e.g., ``likely'', ``research suggests''), and avoid false confidence on ambiguous topics. Penalize overconfident claims and responses that present speculation as fact.

\textbf{9.~Structural Organization and Formatting}\quad
Evaluate the structural organization of the response. Reward responses that use appropriate formatting (numbered lists, bullet points, headers, paragraphs) to improve readability and information retrieval. Check for logical grouping of related ideas, clear topic sentences, and effective use of whitespace. Penalize wall-of-text responses and poorly organized information dumps.

\textbf{10.~Evidence Density and Specificity}\quad
Evaluate the density of concrete evidence and specific details in the response. Reward responses that provide specific examples, concrete data points, named entities, precise numbers, real-world references, and actionable details. Penalize vague, hand-wavy responses that use generic filler like ``many people think'', ``it depends'', or ``there are various factors'' without actually specifying them.
\end{promptbox}

\subsection{Few-Shot Example Format}\label{app:fewshot}

\begin{promptbox}{fewshotbg}{fewshotheader}
\textbf{Few-Shot Example Format (repeated $\times$ 10)}

\medskip
-{}-{}- Example \pgreen{[i]} -{}-{}-

Query: \pgreen{[query text, truncated to 300 chars]}

Response A (score=\pgreen{[score\_a]}):

\pgreen{[response\_a text, truncated to 400 chars]}

Response B (score=\pgreen{[score\_b]}):

\pgreen{[response\_b text, truncated to 400 chars]}

Ground-truth verdict: \pgreen{[Response A is better $|$ Response B is better $|$ Tie]} (gap = \pgreen{[|score\_a $-$ score\_b|]})
\end{promptbox}

\subsection{Avoid-Repetition Block}\label{app:avoid}

\begin{promptbox}{avoidbg}{avoidheader}
\textbf{Avoid-Repetition Block (injected for variant $\geq$ 2)}

\medskip
The following approaches have ALREADY been implemented for this heuristic. You MUST use a substantially DIFFERENT algorithm:

~~-- Variant 1: \pgreen{[approach\_summary]}

~~-- Variant 2: \pgreen{[approach\_summary]}

~~-- ...
\end{promptbox}

\end{document}
